%%%%%%%%%%%%%%%%%%%%%%%%%%%%%%%%%%%%%%%%%%%%%%%%%%%%%%%%%%%%%%%%%%%%%%%%%%%%%%%
\subsubsection{Domain 10.9.3: Financial Decisions}
\label{subsubsec:financial-decisions}
%%%%%%%%%%%%%%%%%%%%%%%%%%%%%%%%%%%%%%%%%%%%%%%%%%%%%%%%%%%%%%%%%%%%%%%%%%%%%%%

\paragraph{The Challenge.}
Household financial decisions---saving, debt management, investment, retirement
planning---show systematic ``biases'' that standard economics cannot explain:
undersaving, over-borrowing, present bias, and failure to diversify. The EBF
framework reveals these as rational responses to multi-dimensional welfare.

\paragraph{Context Profile ($\Psi$).}

\begin{table}[htbp]
\centering
\caption{Financial Decision Context Factors}
\label{tab:finance-psi}
\renewcommand{\arraystretch}{1.3}
\begin{tabular}{@{}lll@{}}
\toprule
\textbf{$\Psi$-Dimension} & \textbf{Effect} & \textbf{Mechanism} \\
\midrule
$\Psi_T$ (Temporal) & Critical & Low stability → short time horizons \\
$\Psi_I$ (Institutional) & Important & Property rights, contract enforcement \\
$\Psi_E$ (Economic) & Direct & Income level constrains options \\
$\Psi_S$ (Social) & Indirect & Trust enables financial cooperation \\
$\Psi_K$ (Information) & Enabling & Financial literacy \\
\bottomrule
\end{tabular}
\end{table}

\paragraph{The Core Trade-off: Consumption Now vs. Security Later.}
Financial decisions involve classic intertemporal choice:

\begin{equation}
\text{Save} \Rightarrow U^{INU}_{F,0} \downarrow \text{ (less consumption now)} \quad \text{but} \quad U^{INU}_{F,3} \uparrow \text{ (more security later)}
\end{equation}

\paragraph{Beyond INU: The Full Picture.}
But financial decisions are not purely individual:

\begin{table}[htbp]
\centering
\caption{Financial Decision: Multi-Category Analysis}
\label{tab:finance-categories}
\renewcommand{\arraystretch}{1.3}
\begin{tabular}{@{}lp{5cm}p{5cm}@{}}
\toprule
\textbf{Category} & \textbf{Saving} & \textbf{Spending/Debt} \\
\midrule
INU & $F_{t=0} \downarrow$, $F_{t=3} \uparrow$ & $F_{t=0} \uparrow$, $F_{t=3} \downarrow$ \\
KNU & Family security ↑ & Immediate family needs met \\
IDN & ``Responsible person'' & ``Generous provider'' / ``Living life'' \\
\bottomrule
\end{tabular}
\end{table}

\paragraph{Reference Points in Financial Decisions.}

\begin{table}[htbp]
\centering
\caption{Financial Reference Points}
\label{tab:finance-reference}
\renewcommand{\arraystretch}{1.3}
\begin{tabular}{@{}llp{5cm}@{}}
\toprule
\textbf{Type} & \textbf{Reference} & \textbf{Effect} \\
\midrule
Status quo & Current consumption level & Saving feels like loss \\
Social comparison & Neighbor's consumption & ``Keeping up with Joneses'' \\
Aspiration & Lifestyle goals & Gap creates dissatisfaction \\
IDN standard & ``People like me have X'' & Identity-driven spending \\
\bottomrule
\end{tabular}
\end{table}

\paragraph{Why People Under-Save: A Welfare Analysis.}

Consider a household deciding whether to save an additional \euro{500/month:

\begin{table}[htbp]
\centering
\caption{Saving Decision: Welfare Components}
\label{tab:saving-welfare}
\small
\renewcommand{\arraystretch}{1.2}
\begin{tabular}{@{}lrrrr@{}}
\toprule
\textbf{Component} & \textbf{Save} & \textbf{Spend} & \textbf{$\Delta$} & \textbf{Discounted} \\
\midrule
\multicolumn{5}{l}{\textit{INU}} \\
$U^{INU}_{F,0}$ (consumption now) & 0.5 & 0.7 & --0.2 & --0.20 \\
$U^{INU}_{F,3}$ (security later) & 0.8 & 0.4 & +0.4 & +0.24 \\
$U^{INU}_{E,0}$ (stress about future) & 0.6 & 0.4 & +0.2 & +0.20 \\
\midrule
\multicolumn{5}{l}{\textit{KNU}} \\
$U^{KNU}_{F,0}$ (family experiences) & 0.4 & 0.7 & --0.3 & --0.30 \\
$U^{KNU}_{F,3}$ (family security) & 0.8 & 0.5 & +0.3 & +0.18 \\
\midrule
\multicolumn{5}{l}{\textit{IDN}} \\
$U^{IDN}_{F,0}$ (``good provider now'') & 0.5 & 0.8 & --0.3 & --0.30 \\
$U^{IDN}_{E,0}$ (``responsible person'') & 0.7 & 0.4 & +0.3 & +0.30 \\
\midrule
\multicolumn{4}{l}{\textbf{Net Welfare Difference (Save - Spend)}} & \textbf{+0.12} \\
\bottomrule
\end{tabular}
\end{table}

The calculation shows saving is \emph{marginally} better (+0.12), but:
\begin{itemize}
\item With slightly higher discounting, it becomes negative
\item With ``good provider'' identity dominant, spending wins
\item With social comparison pressure, spending wins
\end{itemize}

\paragraph{The Debt Trap: Identity and Emotions.}
Consumer debt often serves non-financial functions:

\begin{table}[htbp]
\centering
\caption{Why People Take on Debt}
\label{tab:debt-reasons}
\renewcommand{\arraystretch}{1.3}
\begin{tabular}{@{}lll@{}}
\toprule
\textbf{Utility Target} & \textbf{Mechanism} & \textbf{Example} \\
\midrule
$U^{INU}_E$ & Emotional regulation & Retail therapy \\
$U^{KNU}_S$ & Social belonging & Gifts, experiences with others \\
$U^{IDN}$ & Identity maintenance & ``Person who can afford X'' \\
$U^{INU}_S$ & Status signaling & Visible consumption \\
\bottomrule
\end{tabular}
\end{table}

Debt is not irrational---it trades future $U^{INU}_F$ for present gains 
in other dimensions.

\paragraph{Context Effects on Financial Behavior.}

\begin{table}[htbp]
\centering
\caption{Context Effects on Saving Behavior}
\label{tab:finance-context}
\renewcommand{\arraystretch}{1.3}
\begin{tabular}{@{}lll@{}}
\toprule
\textbf{Context Change} & \textbf{$\Psi$ Effect} & \textbf{Saving Effect} \\
\midrule
Economic instability & $\Psi_T \downarrow$ & $\delta \downarrow$ → Less saving \\
High inequality & $\Psi_S \downarrow$ & Social comparison ↑ → More spending \\
Strong institutions & $\Psi_I \uparrow$ & Trust in future ↑ → More saving \\
Financial literacy & $\Psi_K \uparrow$ & Better understanding → More saving \\
\bottomrule
\end{tabular}
\end{table}

\paragraph{Intervention Strategies.}

\begin{table}[htbp]
\centering
\caption{Financial Behavior Interventions}
\label{tab:finance-interventions}
\renewcommand{\arraystretch}{1.3}
\begin{tabular}{@{}lll@{}}
\toprule
\textbf{Target} & \textbf{Strategy} & \textbf{Example} \\
\midrule
$\delta$ & Reduce discounting & Automatic enrollment, commitment savings \\
$a_{F,t=3}$ & Make future salient & Retirement projections, aging avatars \\
$r^{INU}_F$ & Shift reference point & ``Pay yourself first'' framing \\
$U^{IDN}$ & Identity alignment & ``Savers like you'' messaging \\
$U^{KNU}$ & Social saving & Group saving schemes, tontines \\
\bottomrule
\end{tabular}
\end{table}

\paragraph{Case Example: Automatic Enrollment in Retirement Savings.}
Automatic enrollment works because it changes multiple components:

\begin{itemize}
\item \textbf{Status quo shift:} Saving becomes default → $r^{INU}_F$ adjusts
\item \textbf{Loss aversion:} Opting out feels like giving up employer match
\item \textbf{Effort reduction:} No active decision required
\item \textbf{Social proof:} ``Most people are enrolled''
\end{itemize}

Participation rates increase from ~50\% (opt-in) to ~90\% (opt-out) with 
identical economic incentives.

\paragraph{Key Insight.}
\begin{tcolorbox}[colback=yellow!10!white, colframe=yellow!50!black]
\textbf{Financial Decision Insight:}
``Irrational'' financial behavior reflects:
\begin{enumerate}
\item Rational discounting given uncertainty ($\Psi_T$ low)
\item Non-financial utility from spending (E, S, IDN)
\item Reference point effects (status quo, social comparison)
\item Identity maintenance (``provider'', ``successful person'')
\end{enumerate}
Effective interventions work with these forces, not against them:
automatic enrollment, commitment devices, and identity-aligned messaging.
\end{tcolorbox}

%%%%%%%%%%%%%%%%%%%%%%%%%%%%%%%%%%%%%%%%%%%%%%%%%%%%%%%%%%%%%%%%%%%%%%%%%%%%%%%
% END OF 10.9.3
%%%%%%%%%%%%%%%%%%%%%%%%%%%%%%%%%%%%%%%%%%%%%%%%%%%%%%%%%%%%%%%%%%%%%%%%%%%%%%%
